\section*{Use of AI Tools}
The first author used Claude (Anthropic), a large language model, in the following ways. During the development of the experiments, Claude wrote a large part of the code in the repository. It also translated the repository and the manuscript from Portuguese into English, formatted the manuscript in \LaTeX, checked the references, and proofread the final text. During the revision of the manuscript for submission, Claude, used through Claude Code with the models Claude Fable 5.2 and Claude Opus 5.5, compared the manuscript and the repository with the comments of a pre-submission review and proposed a list of changes. It then reorganized the repository and translated the experimental protocol. It wrote the code that extends the resource counts to larger networks and the code of the additional analyses of the stored data, which it proposed, namely the growth fits, the Wilson intervals, the total variation distances, the frequencies among the shots inside the budget, and the cost of single runs as a fraction of the budget of Proposition~\ref{prop:advantage}. It also redrew Fig.~\ref{fig:example} and drafted the passages of the text that changed in the revision. The first author reviewed and approved all the code and text produced with the tool, and the authors take full responsibility for the content of the paper. The research question, the hypotheses, the design of the experiments, the collection of the data, and the analysis and interpretation of the results obtained before the revision are the authors' own. No AI tool is listed as an author, and no figure was produced by a generative image model.
